% Lecture 11 - Page 1 (13-Dec-1994).
% Fragment for % Lecture 11 - Page 1 (13-Dec-1994).
% Fragment for % Lecture 11 - Page 1 (13-Dec-1994).
% Fragment for % Lecture 11 - Page 1 (13-Dec-1994).
% Fragment for \input{courses/pages/curs11-pagina1-en}.
% No preamble. The first page of a new lecture starts on a fresh page.
% Preserve the manuscript's order; uncertain readings are identified in comments.

\clearpage
\subsection*{Lecture 11}

\begin{center}
% \fbox{\textbf{LECTURE 11}}
\end{center}
\begin{flushright}
\textit{13-Dec-1994}
\end{flushright}

\begin{enumerate}
    \item Sylow's theorems + applications.
    \item Isomorphism theorems + 1 application.
    \item Burnside's theorem + applications.
    \item Table of contents - guide for preparing the algebra examination.
\end{enumerate}

\medskip
\noindent $G$ a group, $|G|=n=p^t m$, $t\geq 1$ and $p\nmid m$.
% In the original, the sign before "p | m" is overwritten/crossed out;
% p\nmid m is the most likely reading and matches the Sylow hypothesis.

\bigskip
\noindent\underline{T$_1$ - Sylow I.}\quad
For every $i$, $1\leq i\leq t$, there exists $H\leq G$ with $|H|=p^i$.

\medskip
\noindent $\operatorname{Syl}_p(G)$ - the set of Sylow $p$-subgroups of $G$,
$\operatorname{Syl}_p(G)\neq\varnothing$.

\bigskip
\noindent\underline{Sylow II.}
\[
G\times\operatorname{Syl}_p(G)\longrightarrow\operatorname{Syl}_p(G),
\qquad (a,P)\longmapsto aPa^{-1}\in\operatorname{Syl}_p(G).
\]
\noindent $\varphi_a(P)=aPa^{-1}$.
% The function symbol written below the mapping appears to be "phi_a(P)".

\medskip
\noindent The action is transitive $\Longleftrightarrow$ it has a single orbit.

\[
\Longleftrightarrow\quad
(\forall)\,P,Q\in\operatorname{Syl}_p(G),\ (\exists)\,a\in G
\text{ such that }aPa^{-1}=Q,
\quad \operatorname{Orb}_G(P)=\operatorname{Syl}_p(G).
\]

\bigskip
\noindent\fbox{Sylow III.}\quad
$|\operatorname{Syl}_p(G)|\mid\operatorname{ord}(G)$ and
\[
|\operatorname{Syl}_p(G)|=pk+1\equiv1\pmod p.
\]

\bigskip
\noindent\fbox{Lemma III.}\quad
Let $P\in\operatorname{Syl}_p(G)$.
\[
\operatorname{Orb}_G(P)=\operatorname{Syl}_p(G).
\]
\[
|\operatorname{Syl}_p(G)|=|\operatorname{Orb}_G(P)|
=[G:\operatorname{Stab}_G(P)]
\overset{\text{Lagrange}}{\Longrightarrow}
|\operatorname{Syl}_p(G)|\mid|G|.
\]

\medskip
\noindent Let $P\in\operatorname{Syl}_p(G)$. Consider the map
(action by conjugation):
\[
P\times\operatorname{Syl}_p(G)\longrightarrow\operatorname{Syl}_p(G),
\qquad (a,Q)\longmapsto aQa^{-1}\in\operatorname{Syl}_p(G).
\]

\noindent Let $A=\operatorname{Syl}_p(G)$.

\medskip
\noindent\underline{Lemma 1.}\quad
$G$ a $p$-group, $|G|=p^a$, $A$ a finite $G$-set,
\[
A_0=\{x\in A\mid ax=x,\ (\forall)\,a\in G\}.
\]
\noindent Then $|A|\equiv|A_0|\pmod p$.

\medskip
\noindent
\[
A_0=\{Q\in A\mid aQa^{-1}=Q,\ (\forall)\,a\in P\}.
\]
\noindent $A_0\neq\varnothing$, since $P\in A_0$
$(aPa^{-1}=P\ \Longleftarrow\ a\in P)$.
% A one-way implication seems to be present in the manuscript.
% The page ends here; the remainder of the argument is not supplied.
.
% No preamble. The first page of a new lecture starts on a fresh page.
% Preserve the manuscript's order; uncertain readings are identified in comments.

\clearpage
\subsection*{Lecture 11}

\begin{center}
% \fbox{\textbf{LECTURE 11}}
\end{center}
\begin{flushright}
\textit{13-Dec-1994}
\end{flushright}

\begin{enumerate}
    \item Sylow's theorems + applications.
    \item Isomorphism theorems + 1 application.
    \item Burnside's theorem + applications.
    \item Table of contents - guide for preparing the algebra examination.
\end{enumerate}

\medskip
\noindent $G$ a group, $|G|=n=p^t m$, $t\geq 1$ and $p\nmid m$.
% In the original, the sign before "p | m" is overwritten/crossed out;
% p\nmid m is the most likely reading and matches the Sylow hypothesis.

\bigskip
\noindent\underline{T$_1$ - Sylow I.}\quad
For every $i$, $1\leq i\leq t$, there exists $H\leq G$ with $|H|=p^i$.

\medskip
\noindent $\operatorname{Syl}_p(G)$ - the set of Sylow $p$-subgroups of $G$,
$\operatorname{Syl}_p(G)\neq\varnothing$.

\bigskip
\noindent\underline{Sylow II.}
\[
G\times\operatorname{Syl}_p(G)\longrightarrow\operatorname{Syl}_p(G),
\qquad (a,P)\longmapsto aPa^{-1}\in\operatorname{Syl}_p(G).
\]
\noindent $\varphi_a(P)=aPa^{-1}$.
% The function symbol written below the mapping appears to be "phi_a(P)".

\medskip
\noindent The action is transitive $\Longleftrightarrow$ it has a single orbit.

\[
\Longleftrightarrow\quad
(\forall)\,P,Q\in\operatorname{Syl}_p(G),\ (\exists)\,a\in G
\text{ such that }aPa^{-1}=Q,
\quad \operatorname{Orb}_G(P)=\operatorname{Syl}_p(G).
\]

\bigskip
\noindent\fbox{Sylow III.}\quad
$|\operatorname{Syl}_p(G)|\mid\operatorname{ord}(G)$ and
\[
|\operatorname{Syl}_p(G)|=pk+1\equiv1\pmod p.
\]

\bigskip
\noindent\fbox{Lemma III.}\quad
Let $P\in\operatorname{Syl}_p(G)$.
\[
\operatorname{Orb}_G(P)=\operatorname{Syl}_p(G).
\]
\[
|\operatorname{Syl}_p(G)|=|\operatorname{Orb}_G(P)|
=[G:\operatorname{Stab}_G(P)]
\overset{\text{Lagrange}}{\Longrightarrow}
|\operatorname{Syl}_p(G)|\mid|G|.
\]

\medskip
\noindent Let $P\in\operatorname{Syl}_p(G)$. Consider the map
(action by conjugation):
\[
P\times\operatorname{Syl}_p(G)\longrightarrow\operatorname{Syl}_p(G),
\qquad (a,Q)\longmapsto aQa^{-1}\in\operatorname{Syl}_p(G).
\]

\noindent Let $A=\operatorname{Syl}_p(G)$.

\medskip
\noindent\underline{Lemma 1.}\quad
$G$ a $p$-group, $|G|=p^a$, $A$ a finite $G$-set,
\[
A_0=\{x\in A\mid ax=x,\ (\forall)\,a\in G\}.
\]
\noindent Then $|A|\equiv|A_0|\pmod p$.

\medskip
\noindent
\[
A_0=\{Q\in A\mid aQa^{-1}=Q,\ (\forall)\,a\in P\}.
\]
\noindent $A_0\neq\varnothing$, since $P\in A_0$
$(aPa^{-1}=P\ \Longleftarrow\ a\in P)$.
% A one-way implication seems to be present in the manuscript.
% The page ends here; the remainder of the argument is not supplied.
.
% No preamble. The first page of a new lecture starts on a fresh page.
% Preserve the manuscript's order; uncertain readings are identified in comments.

\clearpage
\subsection*{Lecture 11}

\begin{center}
% \fbox{\textbf{LECTURE 11}}
\end{center}
\begin{flushright}
\textit{13-Dec-1994}
\end{flushright}

\begin{enumerate}
    \item Sylow's theorems + applications.
    \item Isomorphism theorems + 1 application.
    \item Burnside's theorem + applications.
    \item Table of contents - guide for preparing the algebra examination.
\end{enumerate}

\medskip
\noindent $G$ a group, $|G|=n=p^t m$, $t\geq 1$ and $p\nmid m$.
% In the original, the sign before "p | m" is overwritten/crossed out;
% p\nmid m is the most likely reading and matches the Sylow hypothesis.

\bigskip
\noindent\underline{T$_1$ - Sylow I.}\quad
For every $i$, $1\leq i\leq t$, there exists $H\leq G$ with $|H|=p^i$.

\medskip
\noindent $\operatorname{Syl}_p(G)$ - the set of Sylow $p$-subgroups of $G$,
$\operatorname{Syl}_p(G)\neq\varnothing$.

\bigskip
\noindent\underline{Sylow II.}
\[
G\times\operatorname{Syl}_p(G)\longrightarrow\operatorname{Syl}_p(G),
\qquad (a,P)\longmapsto aPa^{-1}\in\operatorname{Syl}_p(G).
\]
\noindent $\varphi_a(P)=aPa^{-1}$.
% The function symbol written below the mapping appears to be "phi_a(P)".

\medskip
\noindent The action is transitive $\Longleftrightarrow$ it has a single orbit.

\[
\Longleftrightarrow\quad
(\forall)\,P,Q\in\operatorname{Syl}_p(G),\ (\exists)\,a\in G
\text{ such that }aPa^{-1}=Q,
\quad \operatorname{Orb}_G(P)=\operatorname{Syl}_p(G).
\]

\bigskip
\noindent\fbox{Sylow III.}\quad
$|\operatorname{Syl}_p(G)|\mid\operatorname{ord}(G)$ and
\[
|\operatorname{Syl}_p(G)|=pk+1\equiv1\pmod p.
\]

\bigskip
\noindent\fbox{Lemma III.}\quad
Let $P\in\operatorname{Syl}_p(G)$.
\[
\operatorname{Orb}_G(P)=\operatorname{Syl}_p(G).
\]
\[
|\operatorname{Syl}_p(G)|=|\operatorname{Orb}_G(P)|
=[G:\operatorname{Stab}_G(P)]
\overset{\text{Lagrange}}{\Longrightarrow}
|\operatorname{Syl}_p(G)|\mid|G|.
\]

\medskip
\noindent Let $P\in\operatorname{Syl}_p(G)$. Consider the map
(action by conjugation):
\[
P\times\operatorname{Syl}_p(G)\longrightarrow\operatorname{Syl}_p(G),
\qquad (a,Q)\longmapsto aQa^{-1}\in\operatorname{Syl}_p(G).
\]

\noindent Let $A=\operatorname{Syl}_p(G)$.

\medskip
\noindent\underline{Lemma 1.}\quad
$G$ a $p$-group, $|G|=p^a$, $A$ a finite $G$-set,
\[
A_0=\{x\in A\mid ax=x,\ (\forall)\,a\in G\}.
\]
\noindent Then $|A|\equiv|A_0|\pmod p$.

\medskip
\noindent
\[
A_0=\{Q\in A\mid aQa^{-1}=Q,\ (\forall)\,a\in P\}.
\]
\noindent $A_0\neq\varnothing$, since $P\in A_0$
$(aPa^{-1}=P\ \Longleftarrow\ a\in P)$.
% A one-way implication seems to be present in the manuscript.
% The page ends here; the remainder of the argument is not supplied.
.
% No preamble. The first page of a new lecture starts on a fresh page.
% Preserve the manuscript's order; uncertain readings are identified in comments.

\clearpage
\subsection*{Lecture 11}

\begin{center}
% \fbox{\textbf{LECTURE 11}}
\end{center}
\begin{flushright}
\textit{13-Dec-1994}
\end{flushright}

\begin{enumerate}
    \item Sylow's theorems + applications.
    \item Isomorphism theorems + 1 application.
    \item Burnside's theorem + applications.
    \item Table of contents - guide for preparing the algebra examination.
\end{enumerate}

\medskip
\noindent $G$ a group, $|G|=n=p^t m$, $t\geq 1$ and $p\nmid m$.
% In the original, the sign before "p | m" is overwritten/crossed out;
% p\nmid m is the most likely reading and matches the Sylow hypothesis.

\bigskip
\noindent\underline{T$_1$ - Sylow I.}\quad
For every $i$, $1\leq i\leq t$, there exists $H\leq G$ with $|H|=p^i$.

\medskip
\noindent $\operatorname{Syl}_p(G)$ - the set of Sylow $p$-subgroups of $G$,
$\operatorname{Syl}_p(G)\neq\varnothing$.

\bigskip
\noindent\underline{Sylow II.}
\[
G\times\operatorname{Syl}_p(G)\longrightarrow\operatorname{Syl}_p(G),
\qquad (a,P)\longmapsto aPa^{-1}\in\operatorname{Syl}_p(G).
\]
\noindent $\varphi_a(P)=aPa^{-1}$.
% The function symbol written below the mapping appears to be "phi_a(P)".

\medskip
\noindent The action is transitive $\Longleftrightarrow$ it has a single orbit.

\[
\Longleftrightarrow\quad
(\forall)\,P,Q\in\operatorname{Syl}_p(G),\ (\exists)\,a\in G
\text{ such that }aPa^{-1}=Q,
\quad \operatorname{Orb}_G(P)=\operatorname{Syl}_p(G).
\]

\bigskip
\noindent\fbox{Sylow III.}\quad
$|\operatorname{Syl}_p(G)|\mid\operatorname{ord}(G)$ and
\[
|\operatorname{Syl}_p(G)|=pk+1\equiv1\pmod p.
\]

\bigskip
\noindent\fbox{Lemma III.}\quad
Let $P\in\operatorname{Syl}_p(G)$.
\[
\operatorname{Orb}_G(P)=\operatorname{Syl}_p(G).
\]
\[
|\operatorname{Syl}_p(G)|=|\operatorname{Orb}_G(P)|
=[G:\operatorname{Stab}_G(P)]
\overset{\text{Lagrange}}{\Longrightarrow}
|\operatorname{Syl}_p(G)|\mid|G|.
\]

\medskip
\noindent Let $P\in\operatorname{Syl}_p(G)$. Consider the map
(action by conjugation):
\[
P\times\operatorname{Syl}_p(G)\longrightarrow\operatorname{Syl}_p(G),
\qquad (a,Q)\longmapsto aQa^{-1}\in\operatorname{Syl}_p(G).
\]

\noindent Let $A=\operatorname{Syl}_p(G)$.

\medskip
\noindent\underline{Lemma 1.}\quad
$G$ a $p$-group, $|G|=p^a$, $A$ a finite $G$-set,
\[
A_0=\{x\in A\mid ax=x,\ (\forall)\,a\in G\}.
\]
\noindent Then $|A|\equiv|A_0|\pmod p$.

\medskip
\noindent
\[
A_0=\{Q\in A\mid aQa^{-1}=Q,\ (\forall)\,a\in P\}.
\]
\noindent $A_0\neq\varnothing$, since $P\in A_0$
$(aPa^{-1}=P\ \Longleftarrow\ a\in P)$.
% A one-way implication seems to be present in the manuscript.
% The page ends here; the remainder of the argument is not supplied.
